% PDF 143--145；原书第 113--115 页。
\section{答案工程}
\label{sec:ch03-answers}
提示工程帮助我们选择合适的模板。类似地，答案工程（answer engineering）设计和优化答案格式，引导语言模型针对具体的任务或问题，给出尽可能准确且切合上下文的回答。为此，需要考虑不同因素，确定答案形式，再用不同方法搜索答案空间，将答案映射到标签。

\subsection{答案形式}
首先要考虑答案形式（answer shape）。它决定模型输出的粒度，可以细到单个词元，也可以是完整的短语或句子。不同任务需要的答案粒度不同，选好答案形式是基于提示的学习取得成效的关键。答案形式有三种基本类型。

\begin{itemize}
\item \textbf{词元。}从预训练语言模型的词表或其子集中，取一个或多个单独的词元。这种形式常用于情感分类、关系抽取和命名实体识别等分类任务（\refcite{ch03-cui}{Cui et al., 2021}；\refcite{ch03-petroni}{Petroni et al., 2019}；\refcite{ch03-yin}{Yin et al., 2019）}。例如，在情感分类中，模型的答案可以是单个词元，如“positive”“negative”或“neutral”。采用这种形式时，答案空间通常只包含少数几个候选词元，因此属于受约束答案空间。
\item \textbf{片段。}片段（chunk 或 span）是较短的多词元序列，通常与完形填空式提示搭配使用。与词元形式不同，片段的长度不固定，一般属于无约束答案空间。这使它适合问答任务，例如回答“\texttt{Dante was born in [z] .}”这样的提示。
\item \textbf{句子。}按任务的需要，答案可以由一个或多个句子组成，也可以是整篇文档。句子形式的答案通常配合前缀式提示使用，常见于摘要、翻译等需要较详细回答的语言生成任务（\refcite{ch03-radford-2019}{Radford et al., 2019）}。这种形式的答案是无约束的。
\end{itemize}

\subsection{定义答案空间}
用 $Z$ 表示答案空间，即模型针对输入可能给出的所有答案的集合。在许多情况下，这个答案空间映射到一组输出类别，记作 $Y$。答案空间大致分为受约束和无约束两类。

\begin{keypoint}
\begin{itemize}
\item 在受约束答案空间中，输出限于一组预定义的答案。对于输出标签数目有限的任务，如文本分类、实体识别或选择题问答，这种空间很有用。
  \begin{itemize}
  \item 在这种设置中，$Z$ 中的每个元素都映射到标签空间 $Y$ 中的一个元素，作为最终输出。例如，在图~\ref{fig:ch03-mask} 的情感分析任务中，答案“terrible”映射到负面类别。
  \item 这种空间与词元形式的答案相关。
  \end{itemize}
\item 在无约束答案空间中，$Z$ 包含所有词元（\refcite{ch03-petroni}{Petroni et al., 2019）}、定长片段（\refcite{ch03-jiang}{Jiang et al., 2020）}或词元序列（\refcite{ch03-radford-2019}{Radford et al., 2019）}。
  \begin{itemize}
  \item 一般无需把这些输出映射到单独的标签空间，它们本身就是最终输出。
  \item 这种空间与片段及句子形式的答案相关。
  \end{itemize}
\end{itemize}
\end{keypoint}

如上所述，受约束答案空间需要选出组成 $Z$ 的具体词元。以下三节讨论怎样选这些元素。

\subsection{手工建立答案映射}
定义 $Z$ 最简单的办法是手工选择。这需要反复试验。先为各个标签选出不同的答案集合，再用数据检验它们的效果。以情感分析的例子为例，提示为“$x'=$ \texttt{[x] This is a [z] movie.}”。最初取
$Z=\{\text{great}\to\text{positive},\ \text{terrible}\to\text{negative}\}$。
也可以试验
$Z=\{\text{good}\to\text{positive},\ \text{bad}\to\text{negative}\}$。
哪一种表现较好，须由实验确定。

映射不必是一一对应的。多个答案可以对应同一个标签。例如，可以把答案空间扩展为 $Z=\{\text{great},\ \text{terrible},\ \text{good},\ \text{bad}\}$，看概率最高的答案属于正面类（great 或 good），还是负面类（terrible 或 bad），据此选择标签。

\subsection{自动建立答案映射——离散搜索}
有研究者开发了自动选择答案的技术。\refcite{ch03-jiang}{Jiang 等人（2020）}提出一种迭代方法，从初始答案空间 $Z'$ 出发，通过改写扩大它的范围。研究者收集答案—标签对 $(z',y)$，对答案 $z'$ 作回译，找出意义相近、措辞不同的词，再检验这些词的效果。

\refcite{ch03-gao}{Gao 等人（2021）}也提出了一种定义标签词映射 $Z\to Y$ 的方法，目标是使小型验证集上的准确率最大。将填入输入数据的模板交给大语言模型，由模型给出可能的映射，再汇集似然最高的预测。这样，数据集的每个标签都有一份词元排序表，其中的词元是大语言模型预测可用于填补掩码的。可以取前 $n$ 项来建立答案映射。

\subsection{自动建立答案映射——连续搜索}
连续答案搜索直接在模型的嵌入空间中进行。和连续提示一样（见第~3.3.4~节），它解除输出必须映射到已知词元的限制，为进一步优化留出空间。以“positive”和“negative”为例。在自然语言中，这两个词或许最接近所需的输出，但它们的嵌入未必是模型解决该任务时最合适的输出。连续搜索允许模型按更贴近其内部信息表示的方式作答，无需将答案强制转换为自然语言。
